\chapter{Department guidelines}

% \section{c3a}
% \section{cibio}
% \section{cimec}
% \section{cismed}
% \section{dem}
% \section{df}
% \section{dfg}
% \section{dicam}
% \section{dii}
% \section{dipsco}
% \section{disi}
% \section{dlf}
% \section{dm}

\section{Dipartimento di Sociologia e Ricerca Sociale}

All'interno del documento \enquote{\emph{\href{https://infostudenti.unitn.it/it/conseguimento-titolo-corsi-di-laurea-sociologia-e-ricerca-sociale}{linee guida per la redazione dell'elaborato finale dl laurea e della tesi dl laurea magistrale}}}, nella sezione \enquote{consigli redazionali} viene riportato:

\subsection{Formato dell'elaborato finale della laurea triennale}

L'elaborato scritto deve essere lungo all'incirca circa \num{90000} caratteri spazi inclusi (circa \num{13000} parole), carattere Arial o Helvetica; corpo del del testo 12, delle note 10.
La tesi deve essere composta a interlinea \num{1.5} con ampi margini laterali (3cm almeno per ogni lato).

\subsection{Struttura dell'elaborato finale e della tesi magistrale}

Generalmente una tesi di laurea magistrale si compone di una breve introduzione, di quattro/cinque capitoli e delle conclusioni.
La struttura dell'elaborato finale è sostanzialmente identica, anche se non composta da \enquote{capitoli} ma da \enquote{paragrafi}.
\begin{itemize}
    \item nell'\textbf{introduzione} vengono presentati i motivi di interesse della tesi e si delinea il modo in cui verrà svolto il tema;
    \item nel \textbf{capitolo 1} (o, per la laurea triennale, nel paragrafo 1) si passa in rassegna la letteratura rilevante sul tema in questione e vi si inquadra il tema scelto (la letteratura disponibile, gli antefatti storici, i precedenti studi sull'argomento, e cosi via), evidenziando, se del caso, l'originalità del lavoro svolto;
    \item nel \textbf{capitolo 2} (o, per la laurea triennale, nel paragrafo 2) si introduce la metodologia utilizzata;
    \item nei \textbf{due o tre restanti capitoli} (o, per la laurea triennale, nei paragrafi) si presentano i risultati ottenuti dalla propria ricerca (generalmente passando dal generale al particolare, dalle applicazioni più semplici a quelle più sofisticate, dal passato al presente e al futuro).
    \item nella \textbf{conclusione} si presentano i propri giudizi sui risultati ottenuti, sulla rispondenza della metodologia adottata allo scopo della ricerca, sui possibili sviluppi della ricerca.
\end{itemize}

La stesura dell'elaborato finale o della tesi viene svolta sulla base della scaletta condivisa con il/la supervisore o relatore/relatrice cominciando di solito dal Capitolo 1. Non c'è un ordine prefissato, ma è consigliabile che l'introduzione e le considerazioni conclusive vengano redatte per ultime: i loro contenuti sono particolarmente importanti: l'interesse del/della lettore/lettrice tende infatti a concentrarsi in particolare su di esse. Dal momento che riassumono l'intero lavoro, è opportuno
che il resto sia già scritto.

Chiaramente, sia l'elaborato finale che la tesi dovranno anche presentare un frontespizio, un indice iniziale e una bibliografia/sitografia finale.

% \section{Scuola di Studi Internazionali}